\begin{thebibliography}{99}
\addcontentsline{toc}{chapter}{Referencias}
\markboth{Referencias}{}
\bibitem{valencia4}
Miguel Ángel Valencia Bucio. \emph{Análisis Matemático 4, Licenciatura en Física y Matemáticas}. Lecturas 1--23 y 25--28, apuntes del curso proporcionados en la carpeta de trabajo «tesis / análisis / análisis 4», sin fecha indicada. La lectura 19 lleva por error «Lectura 18» en su encabezado. Los archivos originales y su correspondencia se documentan en el registro editorial de esta edición.

\bibitem{reconstruccion24}
\emph{Representaciones de un grupo y de su álgebra de convolución}. Lectura 24: reconstrucción editorial independiente, 18 de octubre de 2026. Apunte de trabajo, con ocho ejercicios y soluciones; no se atribuye a Miguel Ángel Valencia Bucio. Archivo proporcionado en la carpeta del curso: \href{https://drive.google.com/file/d/1S_Am0ooAvf7IKqq7bVcxAlXrou3z1rO0/view}{AnM4-LFM-lect24.pdf}. Su contenido se integra y adapta en el capítulo~\ref{cap:v4-13}.

\bibitem{folland2016}
Gerald B. Folland. \emph{A Course in Abstract Harmonic Analysis}. Segunda edición. CRC Press, 2016. DOI: \href{https://doi.org/10.1201/b19172}{10.1201/b19172}. Referencia para Haar, convolución, representaciones unitarias, dualidad e inversión.

\bibitem{maximenkoConvolucionFinita}
Egor Maximenko. \emph{El teorema de convolución para el grupo cíclico finito $\Z_n$}. Apunte del sitio docente del autor, ESFM--IPN, sin fecha indicada, 4 páginas. \url{https://esfm.egormaximenko.com/analysis/convolution_theorem_for_finite_cyclic_group_es.pdf}. Consultado el 18 de octubre de 2026.

\bibitem{maximenkoPeriodicos}
Egor Maximenko. \emph{Núcleos aproximativos periódicos}. Apunte del sitio docente del autor, ESFM--IPN, sin fecha indicada, 4 páginas. \url{https://esfm.egormaximenko.com/analysis/periodic_approximative_kernels_es.pdf}. Consultado el 18 de octubre de 2026.

\bibitem{maximenkoInversion}
Egor Maximenko. \emph{Inversión de la transformada de Fourier}. Apunte del sitio docente del autor, ESFM--IPN, sin fecha indicada, 4 páginas. \url{https://esfm.egormaximenko.com/analysis/inversion_of_the_Fourier_transform_es.pdf}. Consultado el 18 de octubre de 2026.

\bibitem{lozanoMaximenkoPaleyWiener}
Iris Paola Lozano Vite y Egor Maximenko. \emph{El espacio de Paley--Wiener en la recta real y las series cardinales}. Apunte del sitio docente de Egor Maximenko, ESFM--IPN, 2025--2026, 11 páginas. \url{https://esfm.egormaximenko.com/analysis/Paley_Wiener_space_on_the_real_line_and_cardinal_series_es.pdf}. Consultado el 18 de octubre de 2026.

\bibitem{carleson1966}
Lennart Carleson. «On convergence and growth of partial sums of Fourier series». \emph{Acta Mathematica} \textbf{116} (1966), 135--157. DOI: \href{https://doi.org/10.1007/BF02392815}{10.1007/BF02392815}. Citado exclusivamente para identificar el resultado de convergencia casi en todas partes mencionado como panorama en el capítulo~\ref{cap:v4-04}.
\end{thebibliography}
